\documentclass[preprint,12pt, authoryear, review]{elsarticle}
\usepackage[american]{babel}
\usepackage{lineno, csquotes, tipa, amsmath, caption, afterpage, booktabs}
\usepackage{hyperref}
\modulolinenumbers[5]

\journal{Neuropsychologia}


\begin{document}


\section{Discussion}

The present study investigated how morphological and lexical factors influence visual word recognition and how these effects vary with individual differences in reading skill. By combining behavioral and electrophysiological measures, we sought to clarify when and how readers engage morphological structure during lexical access, and how this engagement depends on \textit{Orthographic Sensitivity} and \textit{Language Proficiency}. The findings converge on a view of morphological processing as an interactive, multi-stage system in which base-level familiarity and family-level connectivity make distinct contributions that are tuned by individual skill.

\subsection{Behavioral Patterns: Additive Facilitation for Words and Morphological Interference for Nonwords}

The reaction time (RT) results demonstrate that morphological and lexical variables exert reliable, dissociable influences on visual word recognition. For words, both \textit{Base Frequency} and \textit{Family Size} facilitated recognition independently, indicating that familiarity with the base morpheme and connectivity within the morphological family each promote faster lexical access. The additive nature of these effects implies that stem-level familiarity and family-level activation operate through partially distinct mechanisms that converge on a shared decision process. 

In contrast, \textit{nonwords} elicited the opposite pattern: morphologically complex pseudowords were rejected more slowly than unstructured controls, and higher base frequency produced further slowing. These findings replicate the morpheme-interference effect and suggest that morphological parsing is at least partly automatic, extending even to illegal combinations. \textit{Family Size} did not influence nonword RTs, implying that network-based facilitation depends on lexical legitimacy. Across both stimulus types, \textit{Orthographic Sensitivity} predicted faster responses overall, reflecting greater fluency in basic decoding, whereas \textit{Language Proficiency} showed no direct RT effect—suggesting that general language ability modulates the \textit{quality} rather than the \textit{speed} of lexical processing.

\subsection{ERP Dynamics: Proficiency-Sensitive Morphological Parsing and Integration}

The ERP results illuminate how morphological information is processed over time. In the \textit{N250} window, which indexes early orthographic–morphological parsing, a three-way interaction among \textit{Family Size}, \textit{Base Frequency}, and \textit{Language Proficiency} emerged for words. High-proficiency readers alone differentiated base frequency as a function of family size, indicating efficient morpho-orthographic access that integrates multiple sources of morphological information early in processing. For nonwords, low-proficiency participants showed stronger N250 complexity effects—particularly when pseudowords were derived from large families—suggesting broader but less selective morphological decomposition. These patterns indicate that Language Proficiency refines early parsing by constraining decomposition to contexts where morphological structure is likely to be meaningful.

In the \textit{N400} window, which reflects semantic integration, proficiency again played a central role. For words, \textit{Family Size} and \textit{Base Frequency} interacted, and this relationship was modulated by proficiency: high-proficiency readers displayed larger and sometimes directionally reversed frequency effects across family sizes, consistent with flexible integration of form and meaning information. For nonwords, a \textit{Complexity} × \textit{Family Size} × \textit{Proficiency} interaction revealed that high-proficiency readers effectively neutralized the usual complexity cost when nonwords were built from bases belonging to large families, whereas low-proficiency readers exhibited exaggerated negativity for the same items. Together, these results show that semantic integration processes are strongly shaped by language experience: skilled readers capitalize on morphological structure to streamline semantic access, while less proficient readers remain influenced by surface form cues even when they are misleading.

\subsection{Individual Differences in Morphological Engagement}

The parallel analyses of \textit{Orthographic Sensitivity} and \textit{Language Proficiency} highlight complementary aspects of individual variability. Orthographic Sensitivity primarily affected global processing speed, consistent with differences in decoding efficiency or attentional control, but had minimal influence on ERP morphology effects. Language Proficiency, in contrast, exerted profound neural modulation at both early (\textit{N250}) and later (\textit{N400}) stages. Proficient readers appeared better able to exploit morphological cues selectively—enhancing facilitation when structural information aligned with lexical knowledge and suppressing it when it conflicted with surface form. This pattern supports models proposing that morphological processing is not fixed but dynamically weighted according to linguistic experience and representational depth.

\subsection{Divergent Roles of Orthographic Sensitivity and Language Proficiency}

The present findings also clarify why Orthographic Sensitivity and Language Proficiency, though both indices of reading skill, exert distinct influences on behavioral and neural measures. These constructs capture complementary but separable dimensions of linguistic expertise. 

\textit{Orthographic Sensitivity} reflects skill in decoding letter patterns and sublexical structure—knowledge of graphemic and morphemic regularities that supports fluent form recognition. This sensitivity is most relevant to early visual-orthographic and morpho-orthographic stages of word recognition, where familiarity with common letter sequences and affix patterns accelerates encoding. In contrast, \textit{Language Proficiency} encompasses broader linguistic competence, including vocabulary breadth, syntactic and semantic knowledge, and the ability to integrate contextual cues. It primarily supports lexical-semantic and integrative processes that operate after initial form decoding.

This distinction helps to explain the divergent behavioral and ERP profiles observed here. Across both words and nonwords, Orthographic Sensitivity consistently shortened RTs, reflecting a general fluency advantage at the level of visual and sublexical processing. High-sensitivity readers were faster because they could efficiently parse plausible morphological units—even in pseudowords lacking meaning. When semantic information is unavailable, as in nonword decisions, these decoding skills remain beneficial, allowing readers to reach a lexical decision quickly based on structural familiarity alone.

By contrast, Language Proficiency produced minimal effects on RTs but robust modulations of ERP components, particularly the \textit{N400}. Lexical decision requires only a superficial familiarity judgment, not deep comprehension; thus, proficiency-linked differences in semantic integration may not translate into measurable response time differences. However, the N400 results revealed that high-proficiency readers engaged more extensive semantic networks, showing stronger and more differentiated responses to interactions among \textit{Base Frequency}, \textit{Family Size}, and \textit{Morphological Complexity}. These patterns indicate richer lexical-semantic representations and greater flexibility in integrating morphological cues with meaning.

High-proficiency readers also demonstrated strategic control over morphological processing. In the ERP data, they modulated their responses depending on whether a stimulus was familiar, meaningful, or morphologically structured. For words, they capitalized on morphological structure to facilitate integration; for nonwords, they suppressed irrelevant activation from familiar-looking affixes. Low-proficiency readers, in contrast, appeared to over-rely on surface form cues, showing exaggerated N400 responses to morphologically complex nonwords that mimicked real derivations.

From a cognitive perspective, Orthographic Sensitivity and Language Proficiency differ not only in \textit{what} they influence but also in \textit{when}. Orthographic Sensitivity dominates early, form-based stages of processing (indexed by the N250), yielding broad speed advantages with relatively mild neural differentiation. Language Proficiency becomes critical later, at the semantic integration stage (indexed by the N400), producing stronger neural modulations but little impact on overt reaction times. This temporal division of labor suggests two partially independent but interacting pathways: one emphasizing rapid, form-level decoding, and another guiding deeper semantic analysis.

To illustrate, consider a morphologically complex nonword such as \textit{refradation}. A reader with high Orthographic Sensitivity might swiftly parse the string into familiar morphemic units (\textit{re-} + \textit{-ation}), recognize that it conforms to legitimate orthographic and morphological patterns, and quickly reject it as a nonword—a fast, form-based decision. A high-proficiency reader, by contrast, would proceed further into semantic evaluation, attempting to access or infer meaning for \textit{refrad}, recognizing the failure of integration, and subsequently suppressing activation—an effortful process reflected in the ERP response rather than in behavioral speed.

In short, Orthographic Sensitivity confers efficiency in early visual and morpho-orthographic decoding, while Language Proficiency governs the depth and flexibility of semantic processing. These complementary abilities jointly scaffold skilled reading: the former ensures fast access to candidate lexical forms, and the latter determines how those forms are interpreted and integrated when ambiguity or morphological complexity arises.


\subsection{Theoretical Implications}

Taken together, the behavioral and ERP data support an \textit{interactive view of morphological processing} in which multiple cues—orthographic, morphological, and semantic—are evaluated in parallel. \textit{Base Frequency} likely supports rapid access to the morphemic stem, whereas \textit{Family Size} contributes through co-activation of semantically and morphologically related forms that aid integration at the semantic level. These cues exert additive effects in overt behavior but interact dynamically in the brain, as revealed by proficiency-linked N250 and N400 modulations. Morphological complexity effects for nonwords further demonstrate that decomposition is an obligatory process, yet its consequences depend on how effectively readers can regulate lexical activation. The differential tuning of these processes by Language Proficiency underscores the role of linguistic experience in shaping the balance between automatic morphological parsing and controlled semantic evaluation.

\subsection{Broader Significance and Future Directions}

By jointly examining RT and ERP measures, this study provides a detailed temporal profile of how morphological familiarity and structural complexity influence word recognition. The results emphasize that proficiency shapes \textit{how} rather than \textit{whether} morphology is engaged: skilled readers integrate morpho-orthographic and semantic information flexibly, whereas less skilled readers rely more heavily on form-based cues. Future work could further dissociate orthographic and semantic contributions by manipulating transparency and by tracking developmental or training-related changes in the timing of N250 and N400 effects. More broadly, these findings highlight that individual differences are not mere statistical noise but a critical source of insight into the adaptive nature of morphological processing in skilled reading.

\end{document}